% Options for packages loaded elsewhere
\PassOptionsToPackage{unicode}{hyperref}
\PassOptionsToPackage{hyphens}{url}
\documentclass[
  ignorenonframetext,
]{beamer}
\newif\ifbibliography
\usepackage{pgfpages}
\setbeamertemplate{caption}[numbered]
\setbeamertemplate{caption label separator}{: }
\setbeamercolor{caption name}{fg=normal text.fg}
\beamertemplatenavigationsymbolsempty
% remove section numbering
\setbeamertemplate{part page}{
  \centering
  \begin{beamercolorbox}[sep=16pt,center]{part title}
    \usebeamerfont{part title}\insertpart\par
  \end{beamercolorbox}
}
\setbeamertemplate{section page}{
  \centering
  \begin{beamercolorbox}[sep=12pt,center]{section title}
    \usebeamerfont{section title}\insertsection\par
  \end{beamercolorbox}
}
\setbeamertemplate{subsection page}{
  \centering
  \begin{beamercolorbox}[sep=8pt,center]{subsection title}
    \usebeamerfont{subsection title}\insertsubsection\par
  \end{beamercolorbox}
}
% Prevent slide breaks in the middle of a paragraph
\widowpenalties 1 10000
\raggedbottom
\AtBeginPart{
  \frame{\partpage}
}
\AtBeginSection{
  \ifbibliography
  \else
    \frame{\sectionpage}
  \fi
}
\AtBeginSubsection{
  \frame{\subsectionpage}
}
\usepackage{iftex}
\ifPDFTeX
  \usepackage[T1]{fontenc}
  \usepackage[utf8]{inputenc}
  \usepackage{textcomp} % provide euro and other symbols
\else % if luatex or xetex
  \usepackage{unicode-math} % this also loads fontspec
  \defaultfontfeatures{Scale=MatchLowercase}
  \defaultfontfeatures[\rmfamily]{Ligatures=TeX,Scale=1}
\fi
\usepackage{lmodern}
\ifPDFTeX\else
  % xetex/luatex font selection
\fi
% Use upquote if available, for straight quotes in verbatim environments
\IfFileExists{upquote.sty}{\usepackage{upquote}}{}
\IfFileExists{microtype.sty}{% use microtype if available
  \usepackage[]{microtype}
  \UseMicrotypeSet[protrusion]{basicmath} % disable protrusion for tt fonts
}{}
\makeatletter
\@ifundefined{KOMAClassName}{% if non-KOMA class
  \IfFileExists{parskip.sty}{%
    \usepackage{parskip}
  }{% else
    \setlength{\parindent}{0pt}
    \setlength{\parskip}{6pt plus 2pt minus 1pt}}
}{% if KOMA class
  \KOMAoptions{parskip=half}}
\makeatother
\setlength{\emergencystretch}{3em} % prevent overfull lines
\providecommand{\tightlist}{%
  \setlength{\itemsep}{0pt}\setlength{\parskip}{0pt}}
\usepackage{bookmark}
\IfFileExists{xurl.sty}{\usepackage{xurl}}{} % add URL line breaks if available
\urlstyle{same}
\hypersetup{
  hidelinks,
  pdfcreator={LaTeX via pandoc}}

\author{\texorpdfstring{}{}}
\date{}

\begin{document}

\section{The Mythopoetics of Global Valuation: Los Against Urizenic
Materialism}\label{the-mythopoetics-of-global-valuation-los-against-urizenic-materialism}

\begin{frame}{Urizenic Materialism and the Generation of ``Allegoric
Riches''}
\protect\phantomsection\label{sec-urizenic-materialism}
To counter the chaotic flux of paper abstraction and the arbitrage
strain of bimetallism, Urizen superimposes geometric limit. In Blake's
mythology, Urizen (often derived from ``Your Reason'' or the Greek
\emph{horizein}, ``to limit'') is the architect of the fallen material
world. Erdman's reading of Urizen as the ``great Work master'' and of
``slave labor'' as war finance, reserves, and commodified human energy
was established in our Introduction {[}@erdman1954prophet{]}. Saree
Makdisi extends this reading, positioning Blake as an anti-capitalist
who rejected the entire framework of commodity exchange, rational
calculation, and empire-driven trade as a ``flattening'' of human
existence {[}@makdisi2003william{]}.

Urizen's mathematically segmented universe mirrors the division of labor
that alienated the 18th-century worker from the product of their hands.
E.P. Thompson extends this analysis via his seminal reading of the poem
\emph{London}, interpreting the ``mind-forg'd manacles'' (line 8; Erdman
p.~70) as literal mechanical mechanisms of market alienation imposed by
the expanding, quantifying logic of industrial capitalism
{[}@thompson1993witness{]}. Thompson's critical intervention
demonstrates how Blake's deliberate imagery systematically indicts the
acquisitive ethic: dividing people, enforcing moral bondage, destroying
joy, and leading to death---all consequences of the reduction of human
life to labor-power exchangeable on the open market. This segmentation
of experience inherently relies upon the materialist reduction of
infinite, qualitative spiritual essence into calculable physical
vessels.

Blake's annotations to Francis Bacon's \emph{Essays Moral, Economical,
and Political} (c.~1798) contain his most concentrated marginalia
against the mercantile worldview. Where Bacon recommended ``the opening
and well balancing of trade; the cherishing of manufactures; the
banishing of idleness'' as prescriptions for national prosperity, Blake
dismissed the entire collection as ``Good advice for Satan's kingdom''
{[}@erdman1988complete{]}. More substantively, in response to Bacon's
axiom that national wealth is necessarily extracted from foreign
loss---``whatsoever is somewhere gotten is somewhere lost''---Blake
annotated: ``Bacon's Philosophy has Ruin'd England.'' This zero-sum
mercantile logic is precisely the Urizenic worldview that undergirds the
bimetallic system: wealth conceived not as qualitative human flourishing
but as a fixed quantum of metal to be captured through competitive
extraction.

Blake interrogates this exact ontological dynamic via the bimetallic
dualism in \emph{The Book of Thel} (1789):

\begin{quote}
``Can Wisdom be put in a silver rod? Or Love in a golden bowl?''
(\emph{The Book of Thel}, 1789) {[}@blake1789thel{]}
\end{quote}

By trapping infinite qualitative attributes (Lunar Wisdom and Solar
Love) into standardized, quantified containers (monetary Silver Rods and
Golden Bowls), the Urizenic state ensures only violent limitation. It is
the poetic expression of Gresham's Law and its underlying arbitrage
boundary: the rod and the bowl are not merely metaphors for human
experience, but precise alchemical designations for the two metals whose
statutory relationship governed the entire British monetary system. The
silver rod of Thel encodes the very silver coinage whose flight from
England---melted and shipped to France and India by the East India
Company---destabilized the fragile bimetallic synthesis. The golden bowl
encodes the hoarded, static gold that Newton's 1717 Mint ratio
artificially elevated above its natural equilibrium. Against this
quantification, Blake arrayed his revolutionary figure of Orc. Often
depicted metaphorically as a pillar of fire or a cosmic serpent, Orc
embodies the raw, unmeasured energy of human passion rising up against
the totalizing theocracy of Urizenic law---a fiery rejection of the
predictability and geometric constraint demanded by the gold standard
bureaucracy.
\end{frame}

\begin{frame}{The Fixing of Labour and the Invention of Allegoric
Riches}
\protect\phantomsection\label{sec-allegoric-riches}
The attempt to fix fluctuating, living subjective value within the
state-mandated algorithmic bounds of the Mint ratio inherently evacuates
the divine meaning of human existence. Blake recognized that the
creation of the modern banking system, the issuance of state debt, and
the fractional generation of credit was an act of misguided alchemy. In
\emph{Jerusalem} (Plate 27, lines 10--14), Blake explicitly targets the
architects of this macroeconomic enclosure:

\begin{quote}
``Shall not the Councellor throw his curb Of Poverty on the laborious?
To fix the price of labour; To invent allegoric riches.'' (Erdman
p.~169) {[}@erdman1988complete; @sklar2011blake{]}
\end{quote}

As Susanne Sklar and modern scholars note, ``allegoric riches'' is a
profound critique of credit creation \emph{ex nihilo}. The Bank of
England held finite gold reserves---by 1797, roughly £1.1 million---but
issued paper credit exceeding £10 million, creating a vast
superstructure of fictional wealth {[}@clapham1944bank{]}. This
``allegoric'' system simultaneously enslaved the laboring class to
compounding debt while enriching a monopolistic financial elite who
controlled the issuance of the currency. Blake's phrase anticipates the
modern critique of central banking's capacity to create money through
lending, generating ``allegoric'' claims upon future production that
bear no necessary relationship to present reality. This materialist
mechanics maps directly onto Blake's ultimate villain: the rigid,
exploitative isolation of the \emph{Selfhood}
{[}@schouten2021lucretius{]}. To accept Urizen's economy is to accept a
world where the living pulse of humanity is subjugated to
quantification.
\end{frame}

\begin{frame}{Entering the Mathematical and Ontological Void of Paper
Credit}
\protect\phantomsection\label{entering-the-mathematical-and-ontological-void-of-paper-credit}
The transition from classical bimetallic friction to the 1797 paper
suspension catalyzed an ontological crisis that Blake internalized into
his epic structures. While gold and silver were idolized as ``Gods of
the Earth,'' paper credit was something even more destabilizing: it was
a pure, systemic void. It substituted verifiable human labor and scarce
elemental reality with infinite, self-referential debt obligations.
\emph{Vala, or The Four Zoas} (c.~1797), drafted as the Restriction
crisis broke, is haunted by this paper deluge. These spectral geometries
evoked new folkloric realities on the streets of London, most notably
the legend of the ``Bank Nun'' (Sarah Whitehead), whose brother Philip
Whitehead was convicted and executed for forging Bank of England notes
during the Restriction years. Driven mad by the state's lethal defense
of its fictional paper, Sarah arrived at Threadneedle Street daily for
over twenty-five years, dressed in black, demanding her brother's
return. For Blake, this intertwining of state execution and fiat forgery
was not merely socio-economic tragedy; it was the literal realization of
his mythological terror---the spectral economies of Urizen plagiarized
into living folk memory {[}@makdisi2003william{]}.

As Ian Baucom illuminates in his study of finance capital and the
Atlantic slave trade, this era of British credit relied on an
epistemology of ``spectral value,'' wherein human life itself became
fully fungible and exchangeable on international ledgers
{[}@baucom2005specters{]}. The speculative instruments of the City of
London---bills of exchange, promissory notes, treasury
bonds---constituted an entirely new semiotic regime in which the
\emph{sign} had been decisively severed from any \emph{referent},
floating freely in recursive financial self-reference without a stable
referent.

Alexander Regier argues that this new financial regime operated via
``fragmentation''---the deliberate shattering of the social contract
into isolated, competing financial claims that could be manipulated by
the state {[}@regier2010fragmentation{]}. If gold was the weight of
Urizen's geometry, paper credit was the abyss of Urizen's chaos. It
created an environment where the ``price of labor'' (Blake's exact
phrase from \emph{Jerusalem} Plate 27) was unmoored from reality,
subject entirely to the inflationary whims of the Bank. For Blake,
returning to metallic standards (the 1816 Coinage Act) was a regression
into tyranny, but accepting unbacked paper without limit was equally
disastrous. The only path forward was transcendence of both quantitative
regimes through imaginative creation---Los's refusal to be enslaved by
another's system (\emph{Jerusalem}, Plate 10) {[}@erdman1988complete{]},
developed in \hyperlink{sec-four-levels-vision}{§ Fourfold Vision}.
\end{frame}

\end{document}
